\documentclass{situationist-preprint}

\title{The Subtraction Argument Reversed}
\msid{MET-01-c}
\msfamily{MET-01}
\msdiscipline{Metaphysics}
\mstype{Alternative proof}
\mskeywords{nothingness; empty world; subtraction argument; principle of indifference; modal measure; natural state}

\begin{document}
\maketitle

\begin{abstract}
MET-01-a argues that the empty world receives measure zero among the possibilities, from a principle of combinatorial plenitude and a principle of combinatorial indifference that gives every combination of potential objects the same measure. Indifference is the premise critics most often reject. We reach the same conclusion without it. The subtraction argument moves from a world with objects to the empty world by removing them one at a time; we run it in reverse, from the empty world to every other world by adding them. A single principle, that adding an object to a world never lowers its measure, then bounds the measure of the empty world by $2^{-n}$ in a space of $n$ potential objects, exactly the value indifference assigns.
\end{abstract}

\section{The premise to be avoided}

The argument of MET-01-a has two premises. Combinatorial plenitude (P1 of MET-01-a) says that for every $n$ there are worlds generated by $n$ potential concrete objects and that every combination of their existence and non-existence is possible. Combinatorial indifference (P2 of MET-01-a) says that each such combination receives the same measure. From the two, MET-01-a (Theorem~4.4) shows that the measure of the empty world in the space $W_n$ of $n$ potential objects is $2^{-n}$, and so less than $1/k$ whenever $n \geq k$.

Plenitude is the premise that the subtraction argument of \citet{Baldwin1996} shares, and few who discuss the question deny it. Indifference is the premise the standard objection targets \citep{Lowe1996,Grunbaum2004}: it is said to be arbitrary, or to presuppose that existence and non-existence are on a par, which is what an opponent of the conclusion denies. MET-01-a defends it. Here we show that it can be replaced by a principle that is strictly weaker and that has an independent motivation, so that the conclusion does not rest on indifference at all.

\section{The reversed argument}

The subtraction argument proceeds from a world $w$ containing finitely many concrete objects, each of which could fail to exist without the existence of any other being required, and removes the objects one at a time until nothing remains \citep{Baldwin1996,RodriguezPereyra1997}. Each step takes a world to a neighbor, a world that differs from it with respect to exactly one object. The reversed argument starts at the other end. From the empty world, every world of $W_n$ is reached by adding, one at a time, the objects that exist in it. Plenitude guarantees that each intermediate world is possible.

We use the notation of MET-01-a (Definition~4.1). A \emph{measure} on $W_n$ is an assignment $\mu$ of a non-negative real number to each world of $W_n$, such that the numbers sum to one. For a world $w$ and a potential object $a$ that does not exist in $w$, $w + a$ is the world that agrees with $w$ except that $a$ exists.

\begin{principle}[Addition]\label{p:add}
For every world $w$ of $W_n$ and every potential object $a$ that does not exist in $w$, $\mu(w) \leq \mu(w + a)$.
\end{principle}

\begin{lemma}\label{l:path}
If exactly $k$ objects exist in a world $w$ of $W_n$, there is a sequence of worlds $e = w_0, w_1, \ldots, w_k = w$, where $e$ is the empty world, each $w_{j+1}$ is $w_j + a$ for some object $a$ that does not exist in $w_j$.
\end{lemma}

\begin{proof}
List the objects existing in $w$ as $b_1, \ldots, b_k$ and let $w_j$ be the world in which exactly $b_1, \ldots, b_j$ exist. Each $w_j$ belongs to $W_n$ by plenitude, $w_0$ is $e$, $w_k$ is $w$, and $b_{j+1}$ does not exist in $w_j$.
\end{proof}

\begin{theorem}\label{t:bound}
If $\mu$ is a measure on $W_n$ that satisfies Addition, then $\mu(e) \leq 2^{-n}$.
\end{theorem}

\begin{proof}
Let $w$ be any world of $W_n$. By Lemma~\ref{l:path} and Addition, applied at each step of the sequence, $\mu(e) = \mu(w_0) \leq \mu(w_1) \leq \cdots \leq \mu(w_k) = \mu(w)$. So $\mu(e) \leq \mu(w)$ for each of the $2^n$ worlds of $W_n$ (MET-01-a, Proposition~4.2). Summing these inequalities over all worlds gives $2^n \mu(e) \leq 1$.
\end{proof}

\begin{corollary}\label{c:vanish}
Under plenitude and Addition, for every $k \geq 1$ and every $n \geq k$, the measure of the empty world in $W_n$ is less than $1/k$.
\end{corollary}

\begin{proof}
By Theorem~\ref{t:bound}, the measure is at most $2^{-n} \leq 2^{-k} < 1/k$.
\end{proof}

Corollary~\ref{c:vanish} is the conclusion of MET-01-a (Theorem~4.4), reached without indifference. The bound of Theorem~\ref{t:bound} is attained exactly when $\mu$ is uniform: equality in the sum requires $\mu(e) = \mu(w)$ for every $w$. Indifference is therefore the limiting case of Addition, the case most favorable to the empty world, and every other measure that respects Addition gives the empty world less.

\section{Why Addition is weaker and why it is true}

The Addition principle is strictly weaker than indifference. Indifference entails it, since equal measures satisfy every inequality of the form $\mu(w) \leq \mu(w + a)$. The converse fails: the measure that gives each world a share proportional to $2^{m}$, where $m$ is the number of objects existing in it, satisfies Addition and is not uniform. A critic who finds indifference arbitrary because it treats existence and non-existence as perfectly symmetric need not accept that symmetry here. She need accept only that non-existence is not favored.

That is the content of the Addition principle: absence is not modally privileged. A world does not occupy a larger share of logical space for lacking an object that another world has. The only ground for denying this would be a view on which non-existence is the natural state of potential objects, so that every existence is a departure requiring a share of space to be paid for. \citet[ch.~2]{Nozick1981} called theories of this kind inegalitarian: they take one state as natural and treat deviations from it as what needs explaining. An inegalitarian theory that takes nothingness as the natural state assumes the answer to the question why there is something rather than nothing, namely that, other things being equal, there would be nothing. It cannot be used as a premise against the measure argument without begging the question. Addition is the denial of that assumption and nothing more.

The principle is also the measure-theoretic counterpart of the third premise of the subtraction argument, that the non-existence of one object never requires the existence of another. That premise says that subtraction never leaves the space of possibilities. Addition says that addition never moves into a smaller part of it. Both express the independence of potential objects that plenitude asserts, and both are accepted by the defenders of the subtraction argument, who take the possibility of the empty world to follow from the modal independence of its would-be inhabitants. As Volume~II of MET-01 is planned to argue, the subtraction argument and its refinements establish the possibility of the empty world without bearing on its measure; the reversed argument shows that the same independence, read as a constraint on measure, bounds that measure from above.

\section{Objections}

\begin{objection}[Leibniz's presumption]
\citet[\S 7]{Leibniz1714} held that nothing is simpler and easier than something. If simplicity confers measure, the empty world, the simplest of all, has more measure than any world that adds an object to it, and Addition fails at its first step.
\end{objection}

\begin{reply}
Simplicity is a ground for preferring theories, and a prior that favors simple hypotheses is a feature of credence, not of the space the hypotheses describe. The measure in question is a measure of the share of logical space a possibility occupies (MET-01-a, Section~3, anticipating Chapter~6 of Volume~I), and the simplicity of a description of a world does not make the world occupy more of that space. Leibniz himself drew from the presumption the conclusion that the existence of something requires a sufficient reason outside the world. The presumption is the inegalitarian assumption of the previous section, and it is not available as a premise in an argument about whether nothingness is privileged.
\end{reply}

\begin{objection}[Indifference in disguise]
Addition compares each world only with its supersets. Any reason to accept those comparisons is a reason to accept the reverse ones, and then indifference follows.
\end{objection}

\begin{reply}
The reverse comparisons, $\mu(w + a) \leq \mu(w)$, would say that existence is not favored over non-existence. We have no need of them, and a theorist who holds that fuller worlds occupy more of logical space, perhaps because they realize more of the combinatorial structure, may reject them consistently with Addition. The asymmetry is principled: the burden is on whoever claims that a state is natural, and Addition denies only the claim that the empty world's state is.
\end{reply}

\section{Conclusion}

The conclusion of MET-01 can be reached without combinatorial indifference. Plenitude, together with the principle that absence is not modally privileged, bounds the measure of the empty world in a space of $n$ potential objects by $2^{-n}$, and the bound tends to zero as the space is enlarged. The argument runs the subtraction argument in reverse, and it shows that the independence of potential objects which makes the empty world possible is also what makes its measure vanish.

\begin{thebibliography}{9}\small
\bibitem[Baldwin(1996)]{Baldwin1996} Baldwin, T. (1996). There might be nothing. \emph{Analysis}, 56(4), 231--238.
\bibitem[Grünbaum(2004)]{Grunbaum2004} Grünbaum, A. (2004). The poverty of theistic cosmology. \emph{The British Journal for the Philosophy of Science}, 55(4), 561--614.
\bibitem[Leibniz(1714)]{Leibniz1714} Leibniz, G.~W. (1714). \emph{Principes de la nature et de la grâce fondés en raison}.
\bibitem[Lowe(1996)]{Lowe1996} Lowe, E.~J. (1996). Why is there anything at all? \emph{Proceedings of the Aristotelian Society, Supplementary Volume}, 70, 111--120.
\bibitem[Nozick(1981)]{Nozick1981} Nozick, R. (1981). \emph{Philosophical Explanations}. Harvard University Press.
\bibitem[Rodriguez-Pereyra(1997)]{RodriguezPereyra1997} Rodriguez-Pereyra, G. (1997). There might be nothing: the subtraction argument improved. \emph{Analysis}, 57(3), 159--166.
\end{thebibliography}

\end{document}
